%!TEX TS-program = pdflatex
\documentclass[11pt,a4paper]{article}

% --- PDFLATEX PREAMBLE BLOCK ---
\usepackage[utf8]{inputenc}
\usepackage[T1]{fontenc}
\usepackage[a4paper, top=2.5cm, bottom=2.5cm, left=2cm, right=2cm]{geometry}
\usepackage[english]{babel}

% Set default font to a clean Sans Serif (Latin Modern)
\usepackage{lmodern}
\renewcommand{\familydefault}{\sfdefault}

% --- Additional Packages ---
\usepackage{amsmath,amssymb}
\usepackage{booktabs}
\usepackage{xcolor}
\usepackage{fancyhdr}
\usepackage{parskip}
\usepackage{float}
\usepackage{placeins} % <--- AJOUT POUR LE FLOATBARRIER

% IMPORTANT OVERLEAF : 
% Retirez l'option "[demo]" ci-dessous lorsque vous compilerez sur Overleaf avec vos vraies images.
\usepackage{graphicx} 

% hyperref must be the last package loaded
\usepackage{hyperref}

% --- Colors ---
\definecolor{darkblue}{HTML}{1a1a2e}
\definecolor{accentblue}{HTML}{0f3460}
\definecolor{lightgray}{HTML}{f5f5f5}

% --- Hyperlinks ---
\hypersetup{
    colorlinks=true,
    linkcolor=accentblue,
    urlcolor=accentblue,
    citecolor=accentblue,
}

% --- Section colors ---
\makeatletter
\renewcommand\section{\@startsection{section}{1}{\z@}{-3.5ex \@plus -1ex \@minus -.2ex}{2.3ex \@plus.2ex}{\normalfont\large\bfseries\color{darkblue}}}
\renewcommand\subsection{\@startsection{subsection}{2}{\z@}{-3.25ex\@plus -1ex \@minus -.2ex}{1.5ex \@plus .2ex}{\normalfont\normalsize\bfseries\color{accentblue}}}
\makeatother

% --- Header/Footer ---
\pagestyle{fancy}
\fancyhf{}
\fancyhead[L]{\small\textcolor{gray}{Commodity Return Forecasting}}
\fancyhead[R]{\small\textcolor{gray}{Nassim Taier}}
\fancyfoot[C]{\small\thepage}
\renewcommand{\headrulewidth}{0.4pt}

\begin{document}

% ============================================================
% TITLE PAGE
% ============================================================
\begin{titlepage}
\centering
\vspace*{4cm}

{\Huge\bfseries\color{darkblue} Commodity Return Forecasting}\\[0.8cm]
{\Large Predicting 1-Week and 1-Month Returns\\Using Machine Learning}\\[2cm]

\textcolor{gray}{\rule{8cm}{0.5pt}}\\[1.5cm]

{\large Nassim Taier}\\[0.4cm]


{\normalsize March 2026}

\vfill
\end{titlepage}

% ============================================================
% TABLE OF CONTENTS
% ============================================================
\tableofcontents
\newpage

% ============================================================
% 1. INTRODUCTION
% ============================================================
\section{Introduction}

We forecast 5-day and 21-day commodity returns using 60 anonymised features. To avoid look-ahead bias, we use a strict walk-forward cross-validation method. We shift target variables by the forecast horizon to prevent overlap, and we leave a gap between training and testing sets. All performance metrics, such as the Information Coefficient (IC) and Sharpe ratio, are calculated on non-overlapping periods. We compare a regularised linear model (ElasticNet) with gradient-boosted trees (XGBoost) and measure how much feature engineering improves the results.

% ============================================================
% 2. DATA
% ============================================================
\section{Data Description}

The dataset contains $T = 6,549$ daily observations from January 2000 to February 2025. The original 60 features are anonymised. Our targets are the forward returns over $h$ days, where $h$ is either 5 or 21.

The data has many missing values, averaging 43.0\%. Since the features have no economic labels, our approach relies entirely on finding statistical patterns rather than using macroeconomic theory.

\begin{table}[htbp]
\centering
\caption{Dataset properties}
\begin{tabular}{@{}ll@{}}
\toprule
\textbf{Parameter} & \textbf{Specification} \\
\midrule
Period & 2000-01-03 to 2025-02-06 \\
Observations ($T$) & 6,549 \\
Raw Features & 60 \\
Prediction Horizons ($h$) & 5 days, 21 days \\
Missing Data Fraction & 43.0\% \\
\bottomrule
\end{tabular}
\end{table}

% ============================================================
% 3. EDA
% ============================================================
\section{Exploratory Data Analysis}

\subsection{Target Distributions and Stationarity}

The return distributions show standard financial characteristics: they are roughly symmetric around zero but have heavy tails (kurtosis of 8.63 for the 1-week horizon and 2.93 for the 1-month horizon). Both return series are strictly stationary based on Augmented Dickey-Fuller tests ($p < 10^{-20}$).

\begin{figure}[htbp]
    \centering
    \includegraphics[width=\textwidth]{figures/eda_distributions.png}
    \caption{Target distributions for the 5-day and 21-day horizons.}
\end{figure}

\subsection{Missing Data Structure}

The dataset exhibits a severe and uneven missing data rate (43.0\% overall). Specifically, 27 out of 60 features have over 50\% missing values. Feature availability increases dramatically in distinct historical steps (e.g., around 2004 and 2013).

This structural missingness poses a significant machine learning risk: because XGBoost natively handles NaNs, the model might learn to use the \textit{presence or absence} of a feature as a proxy for the temporal regime, leading to subtle temporal overfitting.

\begin{figure}[htbp]
    \centering
    \includegraphics[width=\textwidth]{figures/eda_availability.png}
    
    \vspace{0.5cm}
    
    \includegraphics[width=\textwidth]{figures/eda_missing.png}
    \caption{Feature availability over time (top) and missing value percentage per feature (bottom).}
\end{figure}

\subsection{Feature Predictiveness and Stability}

While the median feature-target correlation is weak, the right tail of the correlation distribution reveals exceptionally strong individual predictors (e.g., correlations up to 0.35 with the 21-day return). In standard financial datasets, such extreme individual correlations are rare and warrant caution regarding potential upstream look-ahead bias.

\begin{figure}[htbp]
    \centering
    \includegraphics[width=\textwidth]{figures/eda_target_corr.png}
    \caption{Distribution of Pearson correlations between raw features and targets.}
\end{figure}

Rolling correlations over 252 days show that the predictive power of these features changes significantly over time. We also found dense clusters of highly correlated features ($|r| > 0.9$). This redundancy requires regularization in linear models and favors tree-based models capable of handling collinear inputs.

\begin{figure}[htbp]
    \centering
    \includegraphics[width=0.6\textwidth]{figures/eda_corr_matrix.png}
    \caption{Feature correlation matrix revealing clusters of highly collinear variables.}
\end{figure}

\FloatBarrier

% ============================================================
% 4. FEATURE ENGINEERING
% ============================================================
\section{Feature Engineering}

We expanded the feature set from 60 to 314 using only backward-looking transformations to avoid using future information. These new features are designed to capture recent trends and volatility.

For a feature $x_t$ at time $t$, we calculated rolling z-scores to measure how far the current value is from its recent average:
\begin{equation}
    z_{t}^{(w)} = \frac{x_t - \mu_{t-w:t}}{\sigma_{t-w:t} + \epsilon}
\end{equation}
where $w \in \{21, 63\}$ is the lookback window. Cross-sectional percentile ranks were computed to see how each feature compares to others on the same day, which reduces the impact of extreme outliers.

We measured local stability using 21-day rolling standard deviations for the most correlated features. Feature momentum was calculated as the change over the forecast horizon. Finally, we added lagged versions of the top features, shifted by $h$ days to avoid overlap:
\begin{equation}
    x_{t}^{\text{lag}_k} = x_{t - (h + k)}, \quad k \in \{0, 1, 2, 5, 10\}
\end{equation}

% ============================================================
% 5. MODELING
% ============================================================
\section{Modeling Methodology}

We tested three approaches. First, an ElasticNet regression acts as our linear baseline, using a simple method to fill missing values and scale the data. Second, an XGBoost model is trained on the raw 60 features. XGBoost naturally handles the 43.0\% missing data, so we do not need to fill in missing values manually. Finally, we trained XGBoost on the expanded set of 314 features to see if the engineered features improve performance.

We used 5-fold walk-forward cross-validation. For a given horizon $h$, we stop the training set $h$ days before the test set begins. This gap ensures no future information leaks into the training data.

% ============================================================
% 6. EVALUATION METRICS
% ============================================================
\section{Backtest Design and Risk Metrics}

Portfolio weights $w_t$ are sized based on the strength of the model's prediction, scaled to stay between -1 and 1:
\begin{equation}
    w_t = \text{clip}\left(\frac{\hat{y}_t}{\hat{\sigma}_{\hat{y}}}, -1, 1\right)
\end{equation}
where $\hat{y}_t$ is the prediction and $\hat{\sigma}_{\hat{y}}$ is the standard deviation of recent predictions. We apply a trading cost of 2 basis points for every trade. All metrics (IC, Hit Rate, Sharpe, Max Drawdown) are calculated on non-overlapping periods to provide realistic out-of-sample results.

% ============================================================
% 7. RESULTS
% ============================================================
\section{Results}

The results show that better features lead to higher accuracy (Table \ref{tab:model_comp}). Adding engineered features to XGBoost improves the IC by about 10\% compared to using raw features.

\begin{table}[htbp]
\centering
\caption{Mean IC progression across walk-forward folds}
\label{tab:model_comp}
\begin{tabular}{@{}lcc@{}}
\toprule
\textbf{Model Configuration} & \textbf{IC (5d)} & \textbf{IC (21d)} \\
\midrule
ElasticNet & 0.402 & 0.241 \\
XGBoost (Raw Features) & 0.469 & 0.442 \\
XGBoost (Engineered Features) & 0.517 & 0.492 \\
\bottomrule
\end{tabular}
\end{table}

Table \ref{tab:detailed_metrics} shows the final performance of the engineered XGBoost model. The 5-day model performs very well, achieving a Sharpe ratio of 2.71 and a Calmar ratio of 1.37. The 21-day model also finds a real signal (IC = 0.461), but it suffers from massive volatility, with a maximum drawdown of -113.34\% and a Calmar ratio of just 0.11.

\begin{table}[htbp]
\centering
\caption{Out-of-sample performance metrics (Engineered XGBoost, non-overlapping)}
\label{tab:detailed_metrics}
\begin{tabular}{@{}lcc@{}}
\toprule
\textbf{Metric} & \textbf{5-Day Horizon} & \textbf{21-Day Horizon} \\
\midrule
Pearson IC & 0.505 & 0.461 \\
Hit Rate & 66.9\% & 64.1\% \\
Annualised Sharpe Ratio & 2.71 & 1.51 \\
Maximum Drawdown & -24.39\% & -113.34\% \\
Calmar Ratio & 1.37 & 0.11 \\
Win/Loss Ratio & 2.52 & 2.34 \\
Average Daily Turnover & 40.59\% & 24.50\% \\
$N$ Observations & 765 & 181 \\
\bottomrule
\end{tabular}
\end{table}

\begin{figure}[H]
    \centering
    \includegraphics[width=0.7\textwidth]{figures/pnl_drawdown_1w.png}
    
    \vspace{0.5cm}
    
    \includegraphics[width=0.7\textwidth]{figures/pnl_drawdown_1m.png}
    \caption{Simulated net performance. The 5-day horizon (top) generates compounding returns with contained drawdowns. The 21-day horizon (bottom) experiences severe variance and systemic drawdowns.}
\end{figure}

% ============================================================
% 8. FEATURE IMPORTANCE
% ============================================================
\section{Feature Importance Analysis}

The XGBoost split gains show that engineered features are the most useful. The rolling 63-day and 21-day z-scores provide the most information. While the model uses lagged targets, it relies much more on the normalized features. This shows that knowing the recent volatility context is crucial for making good predictions.

\begin{figure}[htbp]
    \centering
    \includegraphics[width=0.48\textwidth]{figures/feature_importance_1w.png}
    \hfill
    \includegraphics[width=0.48\textwidth]{figures/feature_importance_1m.png}
    \caption{Top 20 features ranked by mean XGBoost split gain for the 5-day (left) and 21-day (right) horizons.}
\end{figure}

\FloatBarrier % <--- LE MUR EST ICI : AUCUNE IMAGE NE PEUT PASSER APRÈS CE POINT

% ============================================================
% 9. LIMITATIONS AND FUTURE WORK
% ============================================================
\section{Limitations and Future Work}

This study only looks at one commodity, so we do not know if these signals work for other assets. Because 43.0\% of the data is missing, the model might be learning from the pattern of missing data rather than real economic trends. Finally, our simple position sizing method is not safe for the 21-day horizon because the drawdowns are too large.

Future work will add stop-losses and volatility targeting to manage risk better. We will also try to optimize the model's settings (like tree depth) and explore models that adapt to different market conditions.

% ============================================================
% 10. CONCLUSION
% ============================================================
\section{Conclusion}

We successfully found predictive signals in short-term commodity returns using machine learning. By using a strict, non-overlapping testing method, our XGBoost model achieved an out-of-sample IC of 0.505 and a Sharpe ratio of 2.71 for the 5-day horizon. While the 21-day model also predicts well (IC = 0.461), trading it without a strict risk management strategy leads to unacceptable losses. Overall, this project shows that transforming features based on recent market conditions significantly improves predictions, but strong risk management is required for slower trading signals.

\end{document}